\thispagestyle{empty}
\begin{center}
{\large \bfseries ABSTRACT}
\end{center}
\vspace{0.4cm}

\noindent
\textbf{Title:} \thesistitle\\
\textbf{Researcher:} \researcher \quad \textbf{Supervisor:} \supervisor\\
\textbf{Department / University:} \department, \institution\\
\textbf{Year of Submission:} \submissionyear

\bigskip
\noindent
The COVID-19 pandemic of 2020--2022 produced a discontinuous shock to
cinema worldwide, but for linguistically and culturally distinctive
markets such as the Tamil film industry the disruption was both
immediate and structural.  Cinema halls in Tamil Nadu were shut for
nearly fifteen months across two lockdown phases; theatrical releases
collapsed from 233 in 2019 to 84 in 2020; OTT acquisition prices for
star-led Tamil films briefly more than doubled before correcting; and
direct-to-OTT releases rose from a handful in 2019 to over 170 by 2025.
This thesis examines, in a single integrated framework, how Tamil-film
viewership transformed and how the industry reconfigured itself in
response.

\medskip
\noindent
The study employs a sequential mixed-methods design.  A structured
questionnaire was administered to a stratified sample of \textbf{400
Tamil film viewers} across Chennai and six other urban / semi-urban
clusters in Tamil Nadu, capturing pre- and post-pandemic theatre
attendance, OTT subscription, genre preference, co-viewing context,
device of consumption, willingness-to-pay and language consumption.
A second instrument, supplemented by semi-structured interviews, was
administered to \textbf{50 industry specialists} -- producers,
directors, OTT acquisition heads, distributors, exhibitors, marketing
professionals, music-label executives and trade analysts -- to capture
shifts in production budgets, theatrical-window negotiations,
acquisition-price indices, star-fee restructuring, marketing-mix
adoption and revenue-mix reconfiguration.  Quantitative analysis used
descriptive statistics, paired $t$-tests, McNemar tests, chi-square
tests of independence, ANOVA and multiple regression in SPSS 28;
qualitative interviews were thematically coded in NVivo 14 using a
constant-comparative approach.

\medskip
\noindent
\textbf{Major findings.}
(i) Mean monthly theatre visits fell by 53\% in the post-pandemic
period (3.1 to 1.5 visits, $p<0.001$), with sharper declines among
35--54 year-olds.  (ii) The mean number of OTT subscriptions per
household more than doubled (1.1 to 2.8) and weekly OTT viewing rose
by 91\%.  (iii) Genre preference shifted decisively from mass-action
and family drama toward thriller, realistic / indie, anthology and
investigative content.  (iv) Star-driven choice fell from 64\% to
46\% while content-driven choice rose to 54\%, a statistically
significant inversion.  (v) The theatrical-to-OTT window shortened
from a 5--8 week norm to 3--4 weeks for the majority of releases.
(vi) Production budgets compressed toward the 5--15 crore band, with
the proportion of greenlit films above 30 crore halving.  (vii) Star
fees corrected most sharply at the A- and B-list tiers ($-22\%$ to
$-34\%$) while the top tier remained relatively insulated.
(viii) The revenue mix of the average Tamil release shifted from
62\% theatrical / 6\% OTT to 44\% theatrical / 28\% OTT.

\medskip
\noindent
The thesis develops a four-quadrant ``Hybrid Equilibrium'' framework
that explains the new structural logic of Tamil cinema -- one in which
theatre and OTT are increasingly complementary rather than
substitutable, content quality has overtaken star power as the
dominant demand driver, and the mid-budget ``story-led'' film has
re-emerged as the central commercial unit.  Theoretical contributions
are made to Diffusion-of-Innovations theory (in modelling OTT adoption
in a non-Western linguistic market), to Uses \& Gratifications theory
(by isolating four post-pandemic gratifications), and to the
disruptive-innovation paradigm (by showing how cultural-industries
incumbents can co-opt rather than be displaced by digital entrants).
Practical recommendations are offered for producers, exhibitors,
streaming platforms, regulators and creative-economy policymakers.

\bigskip
\noindent
\textbf{Keywords:} Tamil cinema; post-COVID viewership; OTT platforms;
streaming; cultural industries; theatrical window; industrial
reconfiguration; Diffusion of Innovations; Uses and Gratifications;
hybrid distribution.
\clearpage
